\section{Budgeted Supervised Meta-Learning}
\label{sec:learning}

The handoff in Figure~\ref{fig:control-handoff} supplies reviewed tasks with selected guidance. The learning process in Figure~\ref{fig:handoff} uses attempts and checks from campaign workspaces to revise that guidance. It repeats within a campaign and carries accumulated lessons into later systems. The heater case shows how a failure can motivate a new rule and define where it applies.

\begin{figure}[t]
\centering
% Supervised rule-evolution overview. Layers group existing records;
% they are not new wiki folders, immutable storage, or qualified skills.
% Stage numbers match handoff.tex; Algorithm 2 gives the detailed procedure.
\begingroup
\hyphenpenalty=10000
\begin{tikzpicture}[
 x=.78cm,y=1cm,font=\sffamily\fontsize{8}{9.4}\selectfont,
 band/.style={draw=black!55,fill=black!4,line width=.7pt,rounded corners=2pt,
   text=black,align=left,text width=11.96cm,minimum height=.64cm,inner sep=5pt},
 evoStep/.style={draw=black!60,line width=.7pt,rounded corners=3pt,text=black,
   align=left,text width=2.58cm,minimum height=3.30cm,inner sep=5pt},
 flow/.style={-{Stealth[length=1.8mm,width=1.3mm]},draw=black!85,line width=.8pt},
 label/.style={font=\sffamily\fontsize{7.3}{8.5}\selectfont,
   fill=white,align=center,inner sep=1.5pt},
 note/.style={font=\sffamily\fontsize{7.3}{8.5}\selectfont,align=left}
]
\node[band,text width=12.65cm,minimum height=.91cm,fill=scpYellowFill,draw=scpRed!70]
 (setup) at (8.7,8.40) {\textbf{Human-approved setup:} goal, scope, assumptions, required checks, access, and time/token limits.\\[2pt]
 \textbf{Inputs:} reviewed sources and a candidate rule, extracted first if needed. Consequential changes require recorded review.};
\node[band,fill=scpBlueFill,draw=scpBlue] (rules) at (8.7,7.35)
 {\textbf{Candidate rulebook}\quad Versions; applicability conditions; required checks; links to predecessor rules.};
\node[band,fill=black!7] (cases) at (8.7,6.45)
 {\textbf{Cases and decisions}\quad Failure explanations; source rationale; reviewed changes; approval and evidence links.};
\node[band,fill=black!3] (evidence) at (8.7,5.55)
 {\textbf{Run evidence}\quad Requirements and artifact versions; formal-check results; fingerprints; recorded resource use.};
\node[evoStep,draw=scpBlue,fill=scpBlueFill] (run) at (2.20,2.75)
 {\parbox[t][2.97cm][t]{2.58cm}{\raggedright\textbf{1\quad Apply and check}\\[4pt]
 Authorized agent applies a rule to an approved task.\\[3pt]
 Generate artifacts; run formal and other checks.}};
\node[evoStep,draw=scpRed!70,fill=scpRedFill] (diagnose) at (6.45,2.75)
 {\parbox[t][2.97cm][t]{2.58cm}{\raggedright\textbf{2\quad Retrieve and explain}\\[4pt]
 Retrieve prior cases and sources.\\[3pt]
 Check the failed attempt's rule, plan, and current evidence.}};
\node[evoStep,draw=orange!75!black,fill=scpYellowFill] (revise) at (10.70,2.75)
 {\parbox[t][2.97cm][t]{2.58cm}{\raggedright\textbf{3\quad Propose rule revision}\\[4pt]
 Revise the procedure, applicability, and required checks.\\[3pt]
 Retain successor draft and predecessor.}};
\node[evoStep,draw=scpGreen!80!black,fill=scpGreenFill] (assess) at (15.00,2.75)
 {\parbox[t][2.97cm][t]{2.58cm}{\raggedright\textbf{4\quad Evaluate and review}\\[4pt]
 Apply candidate. Run formal checks and review requirements.\\[3pt]
 Record required human decisions and outcomes.}};
% Record-layer connections are distinct from the ordered processing steps.
\draw[flow,draw=scpBlue] (rules.west) -- (.18,7.35) -- (.18,2.75) -- (run.west);
\node[label,rotate=90] at (.18,5.90) {selected candidate};
\draw[flow,draw=scpBlue] (run.north east) -- (evidence.south -| run.north east);
\node[label] at (2.20,4.78) {record artifacts,\\checks, and usage};
\begin{scope}[on background layer]
 \draw[flow] (cases.south -| diagnose.north) -- (diagnose.north);
\end{scope}
\node[label] at (6.45,4.78) {retrieve linked\\cases and evidence};
\draw[flow,draw=scpGreen!80!black] (assess.east) -- (17.27,2.75) -- (17.27,7.35) -- (rules.east);
\draw[flow,draw=scpGreen!80!black] (17.27,6.45) -- (cases.east);
\draw[flow,draw=scpGreen!80!black] (17.27,5.55) -- (evidence.east);
\node[label,rotate=90] at (17.27,5.43) {attach results and review};
\draw[flow] (run.east) -- (diagnose.west);
\draw[flow] (diagnose.east) -- (revise.west);
\draw[flow] (revise.east) -- (assess.west);
\draw[flow,draw=scpRed,dashed] (assess.south) -- (15.00,.88) -- (2.20,.88) -- (run.south);
\node[label] at (8.7,.88) {Continue only with approval and remaining budget; retain prior consumption};
\node[band,text width=12.65cm,minimum height=1.0cm,fill=black!4] at (8.7,.10)
 {\textbf{Retained result:} successor draft + predecessor + evidence + review history.\\[2pt]
 \textbf{Task acceptance:} all mandatory checks in the approved scope pass; reviewed contents and matching approval remain current.};
\node[note,text width=13.2cm,anchor=north west] at (.18,-.65)
 {\textbf{Recorded candidates:} the archived heater rules remain drafts awaiting qualification for publication.};
\end{tikzpicture}
\endgroup

\caption{Supervised rule evolution in \workbench{}. A layered presentation inspired by WikiSkill~\cite{tang2026wikiskill} separates recorded evidence, cases and decisions, and candidate guidance. The four stages apply a candidate, interpret prior experience, propose a revision, and evaluate it under review. Updates retain the preceding rule and supporting records. Further attempts require current approval and remaining allowance. Acceptance concerns the reviewed task scope; candidate rules require independent tests before publication for reuse. Section~\ref{sec:handoff} explains the stages; Figure~\ref{fig:evolution} gives a concrete rule history.}
\Description{Three persistent layers store evidence, reviewed cases, and candidate guidance. Four stages apply guidance, interpret experience, propose revised rules, and run formal checks and review within approved scope.}
\label{fig:handoff}
\end{figure}

\subsection{Controlled Handoff}
\label{sec:control-handoff}
A \emph{handoff} supplies an authorized agent with inputs and permitted actions and returns its proposal for host checks. Figure~\ref{fig:control-handoff} follows eight exchanges.

\textbf{1--2. Approve and prepare.} The engineer approves the workflow, scope, access, and limits (1). The harness retrieves current inputs and the selected rule, checks allowance, and specifies permitted actions and output format (2). Evaluator-only artifacts remain separate from dispatched context. Shared task documents and approved reference-derived supplements are declared inputs whose provenance is retained.

\textbf{3--4. Propose and check.} The worker returns an artifact or rule proposal and usage report (3). The host runs formal and other checks and records versions, outcomes, and fingerprints (4).

\textbf{5--7. Retrieve and review.} The engineer retrieves pending requests and evidence (5), records the decision and rationale (6), and the host checks its contents and links it to prior records (7). A request alone grants no authority.

\textbf{8. Continue within limits.} Agents select permitted repairs while approval and allowance remain current. Expanding scope, assumptions, access, or budgets requires renewed human authorization. Prior charges and unresolved outcomes remain recorded~\cite{scpCheckpoint}.

A revised rule selected for another task enters at exchange~2; its outcomes return at exchange~4. Figure~\ref{fig:handoff} explains how this experience informs revision.

\begin{figure}[t]
\centering
% Condensed from public draft-rules.json, GOLDEN-DEFECT-CURATION-001.
% Selected actual revisions draft-1, draft-2, draft-4; draft-3 is omitted.
% Approval of the case premise is distinct from qualification of any rule.
% The input is a documented SMT diagnosis for a contract fragment, not a
% recovered full-system GUMBOX execution trace.
\begin{tikzpicture}[
 x=.78cm,y=1cm,
 font=\fontsize{8}{9.4}\selectfont,
 card/.style={draw=black!65,text=black,rounded corners=2pt,line width=.6pt,
   align=left,inner sep=5pt},
 record/.style={card,text width=3.82cm,minimum height=3.47cm,anchor=north},
 candidate/.style={card,text width=3.82cm,minimum height=3.30cm,anchor=north},
 arrow/.style={-{Stealth[length=1.7mm]},draw=black!75,line width=.7pt}
]
\node[record,draw=scpBlue!75,fill=scpBlueFill!65!white] (requirements) at (2.7,7.25)
 {\begin{minipage}[t][3.12cm][t]{3.82cm}\raggedright\textbf{1. Check reference fragment}\\[2pt]
  Golden reference, NORMAL mode.\\
  $\theta$: temperature; $L,U$: lower and upper\\
  bounds; $h$: heating is on.\\
  \textbf{Assume:} ordering not specified.\\
  \textbf{Guarantee:} $\theta<L\Rightarrow h$;\\
  \hspace*{1.20cm}$\theta>U\Rightarrow\neg h$.\end{minipage}};
\node[record,draw=scpRed!75,fill=scpRedFill!70!white,dashed] (conflict) at (8.65,7.25)
 {\begin{minipage}[t][3.12cm][t]{3.82cm}\raggedright\textbf{2. Retain the conflict}\\[2pt]
  \textbf{Input to the fragment:}\\
  $\theta=98,\ L=99,\ U=97$.\\
  Both conditions hold: heating must be\\
  on and off ($h\land\neg h$).\\
  \textbf{Next action:} review the source.\end{minipage}};
\node[record,draw=scpGreen!75!black,fill=scpGreenFill!70!white] (decision) at (14.6,7.25)
 {\begin{minipage}[t][3.12cm][t]{3.82cm}\raggedright\textbf{3. Review ordering}\\[2pt]
  \textbf{Add:} $L<U$; keep both guarantees.\\
  \textbf{Source:} EA-OI-6 handbook.\\
  \textbf{Decision:} assumption approved.\\
  \textbf{Recheck:} contracts, proofs, tests.\\
  \textbf{Open:} order propagation.\end{minipage}};
\draw[arrow] (requirements.east) -- (conflict.west);
\draw[arrow] (conflict.east) -- (decision.west);
\path[draw=black!50,fill=black!1,rounded corners=2pt,line width=.6pt]
 (-.03,-.52) rectangle (17.33,3.53);
\node[anchor=north west,font=\fontsize{8.2}{9.5}\selectfont\bfseries] at (.12,3.41)
 {4. Revise guidance: retain the failure and require new checks};
\node[candidate,draw=scpBlue!70,fill=scpBlueFill!65!white] (first) at (2.7,2.96)
 {\begin{minipage}[t][2.95cm][t]{3.82cm}\raggedright\textbf{draft-1: Record the correction}\\[3pt]
  \textbf{Change:} require a source citation, approval of the change, and the original failed case.\\[3pt]
  \textbf{Checks:} regenerate contracts;\\
  repeat scoped Verus/GUMBOX checks.\end{minipage}};
\node[candidate,draw=orange!65!black,fill=scpYellowFill!70!white] (second) at (8.65,2.96)
 {\begin{minipage}[t][2.95cm][t]{3.82cm}\raggedright\textbf{draft-2: Clarify responsibility}\\[3pt]
  \textbf{Change:} record the reviewed ordering assumption as an input condition.\\[3pt]
  \textbf{Checks:} retain scoped checks;\\
  track the order across connections.\end{minipage}};
\node[candidate,draw=orange!75!black,fill=orange!10] (third) at (14.6,2.96)
 {\begin{minipage}[t][2.95cm][t]{3.82cm}\raggedright\textbf{draft-4: Check source and build}\\[3pt]
  \textbf{Change:} inspect full source; preserve reviewed assumptions and contracts.\\[3pt]
  \textbf{Checks:} bind proofs to source and build configuration; repeat affected proofs and tests.\end{minipage}};
% Fixed-height cards keep both revision arrows horizontal.
\draw[arrow] (first.east) -- (second.west);
\draw[arrow] (second.east) -- (third.west);
\draw[arrow] (decision.south) -- (14.6,3.53);
\node[card,draw=black!55,fill=black!3,text width=13.0cm,align=center,minimum height=.93cm,inner sep=3pt] at (8.65,-1.18)
 {\textbf{Retained:} source fingerprints link the failure, approved correction, and checks.\\[2pt]
  \textbf{Status:} rules unvalidated (\texttt{ABSTRACT\_NOT\_VALIDATED}); rulebook unpublished (\texttt{DRAFT\_NOT\_PUBLISHED}).};
\end{tikzpicture}

\caption{A heater reference case becomes a source of candidate guidance~\cite{scpCheckpoint}. Steps 1--2 expose conflicting guarantees in the contract fragment. Step 3 records review of an added ordering assumption and the checks required afterwards. Step 4 shows selected rule revisions with links to sources and checks; drafts 1, 2, and 4 are selected from the archive. The added assumption was approved, while the candidate rules await qualification. Horizontal arrows show revision order; intermediate and later drafts are omitted.}
\Description{A temperature contract permits conflicting heater commands without an ordering assumption. Review adds an assumption and records supporting evidence. Three selected draft rules retain that lesson and remain pending independent qualification.}
\label{fig:evolution}
\end{figure}

\subsection{Supervised Rule Evolution}
\label{sec:handoff}
Figure~\ref{fig:handoff} connects run evidence, cases and decisions, and candidate rules through four stages governed by the handoff.

\textbf{1. Apply and check.} Extract an initial candidate from reviewed sources when needed. Apply the selected rule to an approved task, run formal and other checks, and retain versions, outcomes, and resource use.

\textbf{2. Retrieve and explain.} Retrieve the attempt, source context, and review history. Express the reusable structure, procedure, or concept suggested by that experience. Refining a rule requires current failed evidence from that candidate and its approved plan.

\textbf{3. Propose a revision.} State the revised procedure, the tasks it applies to, and the checks it requires. Retain the preceding rule and supporting evidence.

\textbf{4. Evaluate and review.} Apply the candidate with fresh checks and requirements review. Record consequential decisions separately. Acceptance concerns the approved task scope; publishing a reusable rule requires independent qualification.

Missing evidence, approval, or allowance leaves work pending. Continuation retains earlier consumption and remains part of the same workflow. Approved exclusions retain failed outcomes without verification credit~\cite{scpCheckpoint}.




\subsection{From Experience to a Candidate Rule}
Learning turns recorded attempts, diagnostics, and checks into reusable guidance. A \emph{meta-rule} states when it applies, its assumptions, a procedure, and the required checks. The \emph{Generalizer} extracts a procedure shared by related tasks; the \emph{Specializer} supplies the current interfaces, assumptions, and obligations. These strategies use the same model worker. Each proposed rule retains its sources and earlier versions, extending SCP's supervised learning and repair process with persistent review and accounting~\cite{tahat2026scp,scpCheckpoint}.

The file handoff in Section~\ref{sec:execution-environments} makes the reusable knowledge explicit. A new target receives a rule snapshot with its own task inputs. The host archive retains the diagnostic, rule revision, and check that motivated each lesson, so a subsequent attempt can refine the rule and leave the earlier history inspectable. Section~\ref{sec:evaluation} identifies the rules carried between the evaluated cases.

Earlier SCP used a similarity comparison when considering formalizations~\cite{tahat2026scp}. The workbench computes cosine distance to retrieve related rules and to compare contract features after architecture checking. It also records model forecasts with explanations. These signals can help prioritize interpretation and review (Section~\ref{sec:fidelity}).

An early \emph{semantic gate} would use the predicted agreement between English requirements and a proposed interpretation to choose whether to proceed to formalization, refine the interpretation, or request review. Automatic admission and branch ranking are proposed extensions of the available comparison and forecast services. Formal checks and human review then assess the candidate. Appendix~\ref{app:semantic-budget} gives the representations, cosine geometry, and connection to budgeted candidate selection~\cite{foster2026rpm,tahat2026dascslides}.

\subsection{Learning from the Heater Case}
Figure~\ref{fig:evolution} follows a reference correction (upper row) and the candidate guidance derived from it (lower row).

\textbf{(1) Check the reference contract.} Using the notation of Section~\ref{sec:background}, the heater fragment requires $\theta<L\Rightarrow h$ and $\theta>U\Rightarrow\neg h$ in NORMAL mode, without ordering $L$ and $U$.

\textbf{(2) Retain the conflicting input.} The diagnostic script uses $\theta=98$, $L=99$, and $U=97$: both conditions hold, requiring heating on and off simultaneously~\cite{scpCheckpoint}. This exposes the missing constraint in the fragment; $U=97$ lies outside the handbook's full upper-bound range~\cite{lempia2009faa}. Giving one output priority changes a guarantee; filtering tests changes coverage; adding an ordering assumption restricts inputs. Each repair therefore needs an explanation of how it meets the requirement.

\textbf{(3) Review the assumption.} The repair record approves $L<U$ and retains both guarantees. It cites \texttt{EA-OI-6}, the operator-interface environmental assumption in Lempia and Miller's handbook requiring $L\leq U-1^{\circ}\mathrm{F}$~\cite[Appendix A.3.4, p.~A-11]{lempia2009faa}. Strict ordering gives this margin for integer degree settings; other representations require the margin to be stated explicitly. Review records responsibility for the chosen premise; proof and tests assess the repaired artifact, with internal propagation retained as an obligation.

\textbf{(4) Revise the rule.} The figure shows revisions of \texttt{GOLDEN-\allowbreak{}DEFECT-\allowbreak{}CURATION-001}. \texttt{draft-1} requires cited correction, consent, failure retention, regeneration, and checks with Verus and GumboX. \texttt{draft-2} distinguishes environmental assumptions from software guarantees. \texttt{draft-4} adds the complete sources and build context and requires reassessment of regenerated contracts and proofs. Each revision retains its predecessor and supporting evidence; these archived candidates await qualification for publication~\cite{scpCheckpoint}.

Section~\ref{sec:retrieval} explains retrieval of these linked records.

\subsection{Learning Within an Allowance}
Algorithm~\ref{alg:learning} receives the case, reviewed sources, and an optional candidate selected directly or through the retrieval comparison. It can also extract a new draft from the sources. Algorithm~\ref{alg:requirements-fidelity} supplies the development and fidelity-review procedure.
\noindent\emph{Reading the procedures.}
A \emph{budgeted model step} reserves an allowance, calls the model, records usage, and retains the outcome. It returns \textsc{pending} if another check, decision, allowance, or human instruction is needed. A changed input produces \textsc{stale}, which requires renewed review. Every return keeps the campaign history. Section~\ref{sec:evaluation} reports a failure in retaining a detailed call record; the procedures describe the intended order of operations.

\begin{algorithm}[!htbp]
\caption{Learn a candidate rule from retained campaign experience}
\label{alg:learning}
\KwIn{Campaign record $\mathcal{C}$; reviewed sources $M$; optional draft rule $r$.}
\KwOut{Updated $\mathcal{C}$ and $r$; current task outcome $q$.}
\scpPhase{Prepare the rule}
\If{the source or plan is unapproved}{
  \Return \scpStatus{pending} for review\nllabel{ln:learn-review}\;
}
\If{the reviewed source or plan changed}{
  \Return \scpStatus{stale}\;
}
\If{no draft $r$ is supplied}{
  Extract $r$ from $M$ with a budgeted call\;
  \If{no draft is returned}{\Return \scpStatus{pending}\;}
  Retain $r$, its sources, and the call record\nllabel{ln:learn-extract-end}\;
}
\scpPhase{Apply and revise the rule}
\While{recorded permission and budget permit another attempt}{
  Select a currently approved prepared task $T$\nllabel{ln:learn-task}\;
  \If{no such task is available}{\Return \scpStatus{pending}\;}
  Apply $r$ to $T$ in a budgeted call; check; set $q$\;
  \If{$q$ is not \texttt{NOT\_VALIDATED}}{\Return $q$\;}
  \If{the failure uses another rule or plan}{\Return \scpStatus{rejected}\;}
  \If{diagnostic inputs changed}{\Return \scpStatus{stale}\;}
  Revise $r$ from permitted feedback in a budgeted call\;
  \If{no revision is returned}{\Return \scpStatus{pending}\;}
  Retain both rules, failure, and attempts; use the new rule as $r$\nllabel{ln:learn-successor}\;
}
\Return \scpStatus{pending} for more allowance or human guidance\;
\end{algorithm}


The first phase extracts a draft and retains its origin (lines~\ref{ln:learn-review}--\ref{ln:learn-extract-end}). The second applies it to a reviewed task. Rewriting requires current failed evidence from the same rule and plan, recorded as \texttt{NOT\_VALIDATED}; the revised rule keeps links to the failure and earlier rule (lines~\ref{ln:learn-task}--\ref{ln:learn-successor}).

\paragraph{Resource accounting}
\label{sec:resource-accounting}
The campaign's resource ledger retains time, tokens, and attempts across retries. Process deadlines bound execution; token reservations estimate usage, then the ledger records measured consumption and any excess. The estimate may be lower than the usage reported by the provider. Exhausted allowances or unresolved usage stop calls. Approved extensions preserve prior consumption (Section~\ref{sec:interface}).

\paragraph{Qualification for reuse}
Algorithm~\ref{alg:learning} returns a task outcome and a revised candidate. Publishing that rule uses the separate qualification and approval service, which checks evidence for declared examples. The archived heater drafts retain their qualification status~\cite{scpCheckpoint}.
